\begin{frame}[t]{\ev{\tagO\,\tagA}Why siblings may fail together}
\lead{In simulated actomyosin, the Arp2/3 level changes collective behaviour.}
\begin{columns}[T,onlytextwidth]
\column{.28\textwidth}
\begin{center}
\begin{tikzpicture}[x=.8cm,y=.7cm]
\draw[Ink,line width=2pt] (0,0)--(2.5,0);
\draw[Teal,line width=1.5pt] (.55,0)--(1,1.25) (1.4,0)--(1.85,-1.05);
\draw[Teal,line width=1.2pt] (.8,.7)--(.2,1.15) (1.65,-.65)--(2.2,-1.15);
\foreach \x/\y in {.55/0,1.4/0,.8/.7,1.65/-.65}{\fill[Gold] (\x,\y) circle(2.5pt);}
\node[labeltext] at (1.15,-1.65) {schematic, gold = Arp2/3};
\end{tikzpicture}
\end{center}
\column{.68\textwidth}
{\small
\begin{tabularx}{\linewidth}{@{}L{2.3cm}Y@{}}
\toprule
Arp2/3 level & Reported network response\\\midrule
High & stalls\\
Low & contracts\\
Intermediate & loosely connected clusters may collapse suddenly in motor-driven avalanches, similar to experimental ``cytoquakes''\\\bottomrule
\end{tabularx}\par}
\vspace{.1cm}
{\footnotesize We forecast these avalanches with network science and machine learning (Li et al., JPCB 2021).\par}
\end{columns}
\claim{This is motivation only.}
\tagline{Analogy: this does not predict software behaviour.}
\sources{Team results: \href{https://doi.org/10.1073/pnas.1922494117}{Liman et al.\ (PNAS, 2020)}, \href{https://pubs.acs.org/doi/10.1021/acs.jpcb.1c04792}{Li et al.\ (JPCB, 2021)}, \href{https://arxiv.org/abs/2006.06503}{Eliaz et al.\ (PRE, 2020, linker valency)}.}
\note{[0:35] Before test two, we state why siblings are expected to fail together. We simulated branched actomyosin networks and observed this behaviour. At high Arp2/3 levels the networks stall, and at low levels they contract. In between, loose clusters may collapse suddenly, in avalanches. We described the avalanches as an analogy to cytoquakes observed in cells.}
\end{frame}
